\FifthArticle{16}
Als Beispiel diene die Gleichung
\[
x^7+28x^4-480=0
\]
für die
\[
\lambda=\frac{6750}{823543},\quad \log\frac{1}{\lambda}=2{,}0863825
\]
gilt. Die Gleichung hat die erste Form, somit \textit{eine} positive Wurzel, und gehört, da $\lambda$ kleiner als $8$ ist, zum zweiten Fall. Die erste Operation besteht darin, die Gleichung $\log\frac{1}{\lambda}=4A-3B$ zu erfüllen, also, wenn man die Rechnung mit drei Dezimalstellen beginnt, die Gleichung
\[
2{,}086=4A-3B
\]
Schon ein flüchtiger Blick auf die obige Tafel zeigt, dass $A$ zwischen $0{,}5$ und $0{,}6$ zu suchen ist. Es ergibt sich nämlich
\[
\begin{array}{c|c|c}
A&4A-3B&\mbox{Fehler}\\\hline
0{,}5&1{,}643&-0{,}443\\
0{,}6&2{,}109&+0{,}023
\end{array}
\]
woraus sich auf den genaueren Wert $0{,}595$ schließen lässt. Eine neue Rechnung mit den Tafeln mit fünf Dezimalstellen, bei der also $\log\frac{1}{\lambda}=2{,}08638$ zu setzen ist, ergibt
\[
\begin{array}{c|c|c}
A&4A-3B&\mbox{Fehler}\\\hline
0{,}595&2{,}08501&-0{,}00137\\
0{,}596&2{,}08961&+0{,}00323
\end{array}
\]
woraus sich der noch genauere Wert $0{,}5953$ erkennen lässt. Schließlich hat man für sieben Dezimalstellen
\[
\begin{array}{c|c|c}
A&4A-3B&\mbox{Fehler}\\\hline
0{,}5952&2{,}0859279&-0{,}0004546\\
0{,}5953&2{,}0863885&+0{,}0000060
\end{array}
\]
Zum Wert $A=0{,}5953$ muss also noch die Korrektur $-\frac{60}{4606}$ in Einheiten der vierten Dezimalstelle hinzukommen. Ich behalte sie in dieser Form bei, da es, wenn zur Bestimmung von $x$ die erste Formel
\SourceLeaf{101}
\[
x^{7}=\frac{f}{c}
\]
verwendet werden soll, nur darauf ankommt, den entsprechenden Wert von $C$ zu finden. Diesen erhält man, indem man zu dem neben $A=0{,}5953$ stehenden Wert $C=0{,}6935705$ die Korrektur $-\frac{60}{4606}\times798$ hinzufügt, wobei Letztere in Einheiten der siebenten Dezimalstelle zu verstehen ist. Also ergibt sich
\[
\begin{array}{rcl}
C&=&0{,}6935695\\
\log f&=&2{,}6812412\\ \hline
7\log x&=&1{,}9876717\\
\log x&=&0{,}2839531\\
x&=&1{,}9228841
\end{array}
\]

Um die negativen Wurzeln zu finden, schreibt man $x=-y$ und sucht die positiven Wurzeln der Gleichung
\[
y^{7}-28y^{4}+480=0
\]
auf. Diese gehört zur dritten Form. Da $\frac{1}{\lambda}=\frac{823543}{6750}$ größer als $\frac{7^{7}}{3^{3}4^{4}}=\frac{823543}{6912}$, aber kleiner als $2^{7}=128$ ist und zugleich $m$ größer als $n$ ist, gilt der dritte Fall: Es gibt zwei Wurzeln, zu deren Bestimmung die Gleichung
\[
2{,}0863825=3A+7B
\]
erfüllt werden muss. Aus der Schlussbemerkung von Artikel 14 weiß man, dass der eine Wert von $A$ kleiner und der andere größer als $\log\frac{4}{3}=0{,}12494$ sein muss. Auch ergeben sich die Grenzen der Werte von $A$ sofort aus der obigen Tafel mit dreistelligen Logarithmen, nach der man erhält:
\[
\begin{array}{c|c|r}
A&3A+7B&\mbox{Fehler}\\ \hline
0{,}0&2{,}107&+0{,}021\\
0{,}1&2{,}078&-0{,}008\\
0{,}2&2{,}084&-0{,}002\\
0{,}3&2{,}132&+0{,}046
\end{array}
\]
Will man zur genaueren Bestimmung zunächst vierstellige Logarithmen verwenden, so hat man für die erste Lösung anfangs
\SourceLeaf{102}
\[
\begin{array}{r|r|r}
A&3A+7B&\mathrm{Fehler}\\\hline
0{,}05&2{,}0869&+0{,}0005\\
0{,}06&2{,}0847&-0{,}0017
\end{array}
\]
Danach ergeben die fünfstelligen Tafeln
\[
\begin{array}{r|r|r}
0{,}052&2{,}08667&+0{,}00029\\
0{,}053&2{,}08638&0
\end{array}
\]
und schließlich die siebenstelligen
\[
\begin{array}{r|r|r}
0{,}0529&2{,}0863943&+0{,}0000118\\
0{,}0530&2{,}0863660&-0{,}0000165
\end{array}
\]
Hiernach wird
\[
\begin{array}{rcl}
A&=&0{,}0529417\\
\log f&=&2{,}6812412\\
7\log y&=&2{,}7341829\\
\log y&=&0{,}3905976\\
-y&=x=&-2{,}4580892
\end{array}
\]
Für die zweite Lösung sieht die auf ähnliche Weise durchgeführte Rechnung folgendermaßen aus:
\[
\begin{array}{r|r|r}
A&3A+7B&\mathrm{Fehler}\\\hline
0{,}19&2{,}0843&-0{,}0021\\
0{,}20&2{,}0868&+0{,}0004
\end{array}
\]
\[
\begin{array}{r|r|r}
0{,}197&2{,}08627&-0{,}00011\\
0{,}198&2{,}08654&+0{,}00016
\end{array}
\]
\[
\begin{array}{r|r|r}
0{,}1975&2{,}0863805&-0{,}0000020\\
0{,}1976&2{,}0864082&+0{,}0000257
\end{array}
\]
\[
\begin{array}{rcl}
A&=&0{,}1975072\\
\log f&=&2.6812412\\\hline
7\log y&=&2{,}8787484\\
\log y&=&0{,}4112498\\
x&=&-2{,}5778036
\end{array}
\]
\SourceLeaf{103}
Die Gleichung, die uns hier als Beispiel diente, wurde absichtlich so gewählt, dass zwei ihrer Wurzeln nur wenig verschieden sind. In einem solchen Fall sind die Reihen, wie schon oben in Artikel\ 11 bemerkt wurde, wegen ihrer sehr langsamen Konvergenz wenig brauchbar. Auch bei der indirekten Lösung ist zumindest eine schwache Analogie dazu erkennbar: Bei den beiden negativen Wurzeln, die gerade die wenig verschiedenen sind, schreiten die aufeinanderfolgenden Näherungen etwas träger fort als bei der positiven. Ein wesentlicher Unterschied besteht aber darin, dass die sehr langsame Konvergenz der Reihen für sämtliche Wurzeln eintritt, während beim indirekten Verfahren die langsamere Annäherung, die auch nur in geringem Grad spürbar ist, lediglich bei den beiden wenig verschiedenen Wurzeln vorkommt.

